\documentclass[12pt]{article}
\usepackage[english]{babel}
% Replace `letterpaper' with`a4paper' for UK/EU standard size
\usepackage[letterpaper,top=2cm,bottom=2cm,left=3cm,right=3cm,marginparwidth=1.75cm]{geometry}
\usepackage{booktabs}
% Useful packages
\usepackage{amsmath}
\usepackage[colorlinks=true, allcolors=blue]{hyperref}
\usepackage[T1]{fontenc}
\usepackage{babel}
\usepackage{geometry}
\geometry{letterpaper}
\usepackage{graphicx}
\usepackage{verbatim}
\usepackage{afterpage}
\usepackage{amsmath}
\usepackage{amsfonts}
\usepackage{amsthm}
\usepackage{amssymb}
\usepackage{setspace}
\usepackage[round]{natbib}
\usepackage{setspace}
\usepackage{pdfpages}
\usepackage{fullpage}
\usepackage{subcaption}
\usepackage{authblk}
\usepackage{array}
\usepackage{dcolumn}
\usepackage{longtable,lscape}
\usepackage{rotating}
\usepackage{pdflscape}
\usepackage{todonotes}
\usepackage[para]{threeparttable}
\usepackage{tabularx}
\usepackage{marginnote}
\usepackage{xfrac}
\usepackage{ulem}
\usepackage{rotfloat}
\usepackage{float} 
\newtheorem{hyp}{Hypothesis}
\usepackage{footmisc}
\usepackage{float} 
\usepackage{booktabs} 
\usepackage{adjustbox}
\renewcommand{\footnotelayout}{\setstretch{1}}
\usepackage{url}
\usepackage{hyperref}

\bibliographystyle{aer}

\newcommand\Comment[1]{\textcolor{red}{#1}}

\makeatletter
\def\thanks#1{\protected@xdef\@thanks{\@thanks
        \protect\footnotetext{#1}}}
\makeatother

\setlength {\marginparwidth }{2cm}
\begin{document}

\title{
\vspace{-20pt}

Labor Supply under Temporary Wage Increases: Evidence from a Randomized Field Experiment 

\thanks{Ekman: Karlstad University, Karlstad Business School, 651 88 Karlstad, Sweden (e-mail: mats.ekman@kau.se). 
Jakobsson: Karlstad University, Karlstad Business School, 651 88 Karlstad, Sweden, and FBK-IRVAPP, Trento, Italy (e-mail: niklas.jakobsson@kau.se). 
Kotsadam: Ragnar Frisch Centre for Economic Research, Gaustadalleen 21, 0349 Oslo, Norway (e-mail: andreas.kotsadam@frisch.uio.no).  We are grateful to the staff at \textit{Situation Sthlm} for their smooth and efficient collaboration, and gratefully acknowledge the financial assistance of the Jan Wallander and Tom Hedelius Foundation, and Tore Browald’s Foundation (Grant no. P23-0086). Refine.ink was used to check the paper for consistency and clarity. An analysis plan \citep{preanalysisplan} is pre-registered at The American Economic Association’s registry for randomized controlled trials (No. AEARCTR-0014308), and all deviations from the plan are noted in the text. The plan is also added to the Appendix.}
}

\author{Mats Ekman, Niklas Jakobsson, and Andreas Kotsadam}

\maketitle
\vspace{-2em}  % adjust -2em as needed (try -1em, -1.5em, etc.)
\thispagestyle{empty}

\begin{abstract}
\noindent 
We conduct a pre-registered randomized controlled trial to test for income targeting in labor supply decisions among sellers of a Swedish street paper. These workers face liquidity constraints, high income volatility, and discretion over hours. Treated individuals received a 25 percent bonus per copy sold for the duration of an issue, simulating an increase in earnings potential. Treated sellers sold more papers, worked longer hours, and took fewer days off. These findings contrast with studies on intertemporal labor supply that find small substitution effects. Notably, when we apply strategies similar to observational studies, we recover patterns consistent with income targeting.

\small

\vspace{2ex}

\noindent \textsc{Keywords}: Field Experiment, Labor Supply, Reference Dependence, Street Paper
\newline

\noindent \textsc{JEL codes}: J22, C93, D91, D03
\newline

\end{abstract}

\clearpage
\doublespacing

\pagebreak

\section{Introduction}
The decision to work less when paid more is consistent with rational behavior due to the possibility of a backward-bending labor supply curve. However, if the higher pay is transitory, standard economic theory predicts that individuals should work more, taking advantage of the temporary incentive to earn more now and enjoy greater leisure in the future. The advent of prospect theory \citep{Kahneman1979prospect} has motivated models with reference-dependent preferences that depart from the standard life-cycle framework, in which consumption smoothing is an implication \citep[e.g.,][]{friedman1957theory, modigliani1954utility}. In prospect theory-inspired applications to labor supply, effort becomes more costly near certain reference points, leading individuals to work longer when earnings are low and reduce effort once a target is met. This type of reference dependence implies that raising wages could perversely reduce labor supply, a concern for both firms and policymakers.

Because of this, there has been growing interest in studying workers who face uncertain incomes and set their own hours. Pioneering work by \citet{camerer1997labor} found strong evidence of reference-dependent behavior among New York City cabdrivers, using observational data. Since then, the taxi industry has been studied further by \citet{farber2005tomorrow, farber2008reference, farber2015taxi}, \citet{crawford2011newyork}, \citet{agarwal2015singaporean}, \citet{martin2017when}, \citet{thakral2021daily}, and \citet{duong2023taxi}, among others. Some of these findings challenge Camerer \textit{et al}.’s conclusions, while others attempt to clarify the conditions under which income targeting is most likely to occur, such as particular income thresholds or planning horizons.

We are aware of only two randomized controlled trials that explicitly address intertemporal substitution of labor supply \citep{fehrgoette,andersen2025toward}. In contrast to some of the taxi studies, they indicate positive labor supply elasticities under conditions of high but transitory earnings potential, although they differ in how these results are interpreted. \cite{fehrgoette} increased the earnings potential of Swiss bike messengers and found a positive labor-supply response but a negative elasticity of effort per hour. Additional lotteries make them conclude that reference-dependence fits the data best. In a two-pronged setup, \cite{andersen2025toward} (1) use confederates to bargain with vendors in an Indian open-air market, such that vendors earned more when the confederates bargained less skillfully, and (2) observed work hours by vendors afterwards, finding that treated vendors worked longer hours, although their response came with a one-day lag.

Our approach differs from most previous analyses in at least three key methodological respects. Firstly, our study is a field experiment where we randomly assign sellers of a major Swedish street paper, \textit{Situation Sthlm} (short for ‘Stockholm’), into treated and control sellers. Street papers are sold by homeless or hard-up individuals, typically as a way for them to come into gainful employment.\footnote{The first street paper was the Chicago-based \textit{Hobo News} in 1915, and current examples include \textit{The Big Issue} (the UK), \textit{StreetWise} (Chicago), and \textit{StreetSheet} (San Francisco).} The sellers are paid solely on commission. To mimic lucky business conditions, treated sellers receive a bonus of ten SEK (about one US dollar) per copy sold. Sellers otherwise receive 40 SEK per copy, so our intervention is a 25 per cent bonus. We pay the bonus on sales of the October issue of the street paper, from the day of its release until the release of the November issue.

Treated sellers increased both their sales and working hours while taking fewer days off. We find no evidence for income-targeting behaviour of the type commonly inferred in observational studies. Yet, when we analyze non-experimental sales data, we reproduce patterns suggestive of income targeting—demonstrating how experimental identification is crucial for causal interpretation.

Secondly, following \citet{leeson2022hobo}, it is particularly worthwhile to examine behavior in persons who are frequently (rightly or wrongly) thought to behave less rationally. \citet{leeson2022hobo} study where panhandlers decide to locate along the Washington DC Metro System, finding that their time-adjusted earnings tend towards equality (partly reproduced by \citet{demeloetalreplication}). Similarly, if we find positive labor supply elasticities over transitory income shocks in sellers of street papers, it is stronger evidence of rational behavior in general than if we had studied members of some other occupation where workers are thought to behave more rationally.

Thirdly, our measurement also adds perspective to prior literature on the elasticity of labor supply. Related to the issue of division bias \citep{borjas1980wages}, a risk in some of the studies that have focused on taxi drivers is that the tendency to use trip-level data for individual drivers may confuse leisure and work. Once on the road, a leased taxi is a sunk cost, and so what may look like work hours may actually be drivers running personal errands, making work hours on slow days appear longer than they actually are. Some studies, such as \citet{duong2023taxi}, use Uber and Lyft to address the issue of when drivers actually work. Services like Uber and Lyft provide data for when their drivers “work” in the sense of waiting for customers, as opposed to when they are unavailable (because with customers or because not working), but drivers may nevertheless be available while running personal errands and so not actively engaged in working. Because our outcome is sales during the period when an issue is current and the bonus is received, we avoid this potential measurement error by relying on an administratively recorded output measure rather than constructing hours from noisy activity/availability data. Under the monotonicity assumption that total sales increase when work time increases, sales provide our main proxy for time spent working for both a control group and a treated group whose members effectively receive higher pay while working. More detailed measures of input—Hours and Days based on electronic timestamps—are introduced later and used as complementary, approximate outcomes.

The remainder of this paper proceeds as follows. In Section 2, we present the experimental design and the data set. In section 3, we present the pre-specified empirical strategy, our findings, heterogeneity analyses, a series of robustness checks, some survey evidence, and an analysis of observational data. Section 4 concludes.


\section{Experimental Design and Data}
\subsection{Experimental Design}
\textit{Situation Sthlm} was founded in 1995 and has an annual circulation of approximately 200,000 copies, sold by about 250 sellers. The sellers are a heterogeneous group, but tend to have no permanent place of residency, and are more likely than most to be troubled by various mental problems and drug addictions. Apart from making money by selling the street paper, they make a living on social assistance, mendicancy, and collecting empty bottles and cans for the deposit. Some sellers occasionally write articles or do publicity work for the street paper. When they sell the paper, they must first purchase it from the office of \textit{Situation Sthlm}, located in Central Stockholm. At the time of the intervention, the price of a copy was 40 SEK to sellers, and the resale price was 80 SEK. The sellers have designated spots across Stockholm where they are encouraged to sell, but they are not restricted to these locations. Unsold copies must be returned by the release of a new issue, and \textit{Situation Sthlm} refunds the entire 40 SEK per copy in such cases. A new issue is released on the last Wednesday of every month except for July (the June issue is a double issue). The sellers may come and go to the office as they please during the office’s hours of operation.

When an electronic sale is made, the payment system is such that the seller can obtain the money from the office of \textit{Situation Sthlm} after about 45 minutes. For severely budget-constrained sellers, this fact implies a cap on daily sales, which is not necessarily raised by our bonus. For example, a seller who has only a hundred SEK could buy two copies of the paper (80 SEK) during the intervention, and come back to the office to obtain 180 SEK with our bonus (enough for five new copies if the seller retained the 20 SEK left over from the first transaction) when these sales have been made and the 45-minute processing of the payment has been concluded. If the seller cannot sell these five papers by the time the office closes, the bonus will increase sales for this individual by just one copy on this day. During the intervention, the hours of operation of the office were from 8 am to 3.15 pm, Monday to Friday, and from 10 am to 3 pm on Saturdays.

On September 23rd, 2024, an announcement was placed on a wall in the office lobby of \textit{Situation Sthlm}, informing readers that the October issue of \textit{Situation Sthlm}, to be released on September 25th, would come with a “Swish campaign”. Swish is the largest mobile payment system in Sweden, and approximately 80 percent of all copies of \textit{Situation Sthlm} are sold using it. The announcement also said that those sellers who wish to participate in the campaign will watch a coin being flipped, and if it comes up heads, they receive ten SEK extra per copy sold on Swish for the October issue, released September 25th. If the coin comes up tails, they instead receive a bonus for the next issue, to be released on October 30th (called the November issue). Hence, the assignment followed a randomized waitlist-control design, with delayed treatment. Sellers were thus assigned a group as they came into the office and agreed to participate. A total of 53 sellers were assigned to the October treatment, and 56 sellers to the November treatment. Five sellers chose not to participate.

Tying the bonus only to Swish sales was done to prevent sellers from purchasing papers from the office for immediate resale to colleagues randomized not to get the bonus\footnote{This is also the reason for the departure from letting treated sellers purchase copies at a discount, mentioned in our Pre-analysis plan (see section A.3 in the appendix). The bonus comes to the same thing, but avoids the risk of resale.}. The paper has a pre-existing system to prevent fraud, in which the sellers must consent in writing not to use the payment system for private reasons. Essentially, the office has access to the names and phone numbers of all customers who pay using Swish, as well as the amounts received, so any time one customer buys several copies or any time the amount paid differs from 80 SEK, the seller knows that he or she will have to explain this to the office. \textit{Situation Sthlm} incurs a cost to enable customers to pay using Swish, so, naturally, should not wish to pay for non-sale transactions.


\subsection{Dataset}
\textit{Data sources:} We use data from the editorial office of \textit{Situation Sthlm} on all sales by individual sellers. For electronic sales, we have access to real-time sales data; for cash payments, we know that the issue was sold between when the seller picked it up and when they went back to return any unsold magazines. In addition to the sales data, we also have access to the sellers' ages and sex. All \textit{Situation Sthlm} sellers in Stockholm who entered the \textit{Situation Sthlm} office during the intervention period got the offer to participate in the experiment. 109 sellers agreed to participate, and five declined. A member of staff flipped a coin in front of the participant, and in this manner, 53 sellers were randomized to the treatment group and 56 to the control group. During the study period, there were eight active sellers not located in Stockholm and therefore not directly in touch with the \textit{Situation Sthlm} office. They were not asked to participate.

\textit{Variable definitions:} Our main outcome variable is sales of the October issue (Sales), which includes all sales made of the issue. Our pre-analysis plan also includes the natural logarithm (+1) of the sales variable (ln(Sales)) and sales via the electronic system only (Electronic), which excludes cash sales, as exploratory outcome variables. Such outcomes also include cash sales only (Cash), which are not-directly-observed cash transactions but rather a residual constructed by subtracting Electronic sales from Sales, where Sales is inferred from inventory movements and Electronic is measured from Swish transactions. Cash can therefore occasionally take negative values. In our data, Cash has eight negative observations, mainly because some sellers may accept Swish payments without handing over a magazine, or because some buyers pay for several issues while only taking one physical copy, so that Swish payments can exceed the number of copies leaving inventory. We also examine time spent working as exploratory outcome variables. Hours worked during the month are approximated by the the span between the first and last electronic sale each day (Hours). The number of working days is also approximated by the total number of days in which a seller sells an issue via the electronic system (Days). In common with some of the taxi studies on supply elasticity \citep[e.g.,][]{camerer1997labor,farber2005tomorrow,crawford2011newyork}, we cannot tell whether a day without sales is a working day (an unsuccessful one) or not. We use binary variables to denote these days. Lastly, the control variables for Age (Age) and sex (Woman) are reported to us by \textit{Situation Sthlm}, and are based on the personal identity numbers of the sellers (these are the Swedish national identification numbers, which are structured to reveal their holders’ sex and year of birth).

\textit{Summary statistics:} Appendix Table \ref{tab:summary} presents summary statistics for the key variables in the study. The average age of the participants is approximately 55.8 years (SD = 10.87), with a minimum of 26 years and a maximum of 82 years. About 34 percent of the sample are women. The mean number of magazines sold during the October issue is 112.85 (SD = 191.07), with considerable variation, ranging from 0 to 1,555. The natural logarithm of sales (ln(Sales)) is also reported, with an average of 3.61 (SD = 1.79). Sales are further disaggregated into electronic and cash sales. Electronic sales constitute the majority, with a mean of 96.36 (SD = 177.08) and a maximum of 1,467. Cash sales are lower on average, at 16.50 (SD = 25.38), with values ranging from -10 to 160. As indicated above, Cash sales include a few negative observations, mainly because some sellers may sell a magazine to a customer without actually handing it over to the buyer, and some buyers may also buy several magazines while only taking one. The participants worked an average of 30 hours per month (SD = 42.68), with some individuals not working at all (i.e., only making cash sales) and others working up to 286.7 hours. The number of days worked during the October issue averages 11.17 (SD = 9.55), with a range from 0 to 33 days.

\textit{Randomization tests:} To assess the validity of the randomization, we test for balance in the control variables between treated and controls. Appendix Table \ref{tab:balance} presents baseline means for treated (n = 53) and controls (n = 56), their differences, and p-values after t-tests. Few differences are statistically significant, but sales in the control group have been consistently higher in the year before treatment. As a joint test of balance, we regress the treatment dummy on the control variables indicating imbalance at baseline (F = 7.16, p = 0.00), implying that the baseline variables jointly predict treatment status. Given the small sample size, some imbalance is not unexpected; we therefore control for pre-treatment sales behavior in all analyses. Including lagged outcomes also improves precision, as prior values are strongly correlated with the outcome (see Appendix Table \ref{tab:correlations}).


\section{Results}
\subsection{Pre-specified regression model}
We estimate several versions of equation~\eqref{eq:regression}:

\begin{equation}
    Y_i = \alpha + \beta_1 T_i + \beta_2 X_i + \varepsilon_i,
    \label{eq:regression}
\end{equation}

where $Y_i$ is an outcome for seller $i$, $T_i$ is an indicator of treatment, and $X_i$ is a set of baseline controls, including age, sex, and 12 lags of the outcome variable. The regressions are estimated with heteroscedasticity-robust standard errors.

We use the post-double LASSO selection approach by \citet{belloni2014inference}. The LASSO selection approach selects variables correlated with treatment and outcomes, improving the precision of the estimates and helping correct imbalances across groups. This is especially important since sales in the treated group were consistently lower in the year leading up to treatment.

\subsection{Main findings}
In Table \ref{tab:main}, we present results for the overall population of sellers using equation~\eqref{eq:regression} and the post-double LASSO selection approach. The coefficient of our main outcome variable, Sales, is 23.25, and the mean sales in the control group is 110.68, implying a 21 percent increase in issues sold following treatment. For ln(Sales), the effect is bigger but less precisely estimated. Judging from the coefficients of Electronic sales and Cash sales, the increase is driven by electronic sales, as expected. The coefficients on Hours worked were 10.99 (control group mean 27.11), and on working Days 2.56 (control group mean 10.66), implying large effects on working time. In this case, full monthly reference dependence would mean a decrease of 20 percent in the number of issues sold, since the price increase was 25 percent. Instead, we find evidence of sales and sales efforts increasing.

\begin{table}[H]
\centering
\begin{threeparttable}
\caption{\label{tab:main} Main results}
\normalsize   % <-- gör texten lika stor som i de andra tabellerna
\input{results/main}
\begin{tablenotes}[flushleft]
\small
\item \textit{Notes:} This table reports estimates of treatment effects for sales and five additional variables that capture sales or working time. Estimations were made using the post-double LASSO selection approach with a full set of potential control variables. Robust standard errors are in parentheses. P-values are $\leq 0.01^{***}$,  $\leq 0.05^{**}$, and $\leq 0.1^{*}$.
\end{tablenotes}
\end{threeparttable}
\end{table}


\subsection{Heterogeneity}
Following previous research \citep[e.g.,][]{farber2015taxi}, there is evidence that income targeting is stronger among inexperienced individuals and weaker among those with more experience. Table \ref{tab:hetero} presents results from heterogeneity analyses based on sellers’ experience levels and age groups. Experience is measured in three ways: (1) by the number of magazines they sold during the last month (standardized with a mean of zero and a standard deviation of one), (2) by the number of magazines sold during the last year (standardized with a mean of zero and a standard deviation of one), and (3) by a dummy variable including sellers that sold more magazines than the median seller during the past three months, as described in the pre-analysis plan. As before, we use the post-double LASSO selection approach with a full set of potential control variables.

\begin{table}
\caption{\label{tab:hetero} Heterogeneity}
\vspace{-0.8em}
\begin{threeparttable}[H]

\begin{subtable}[h]{\textwidth}
\caption{\label{tab:h_a} Sales last month (std)}
\resizebox{\textwidth}{!}{\scriptsize \input{results/h_a}}
\end{subtable}

\begin{subtable}[h]{\textwidth}
\caption{\label{tab:h_b} Sales last year (std)}
\resizebox{\textwidth}{!}{\scriptsize \input{results/h_b}}
\end{subtable}

\begin{subtable}[h]{\textwidth}
\caption{\label{tab:h_c} Sales last three months dummy (median 467)}
\resizebox{\textwidth}{!}{\scriptsize \input{results/h_c}}
\end{subtable}

\begin{subtable}[h]{\textwidth}
\caption{\label{tab:h_d} Age (std)}
\resizebox{\textwidth}{!}{\scriptsize \input{results/h_d}}
\end{subtable}

\begin{tablenotes}
\small
\begin{minipage}{\textwidth}
\textit{Notes:} Each column corresponds to a different outcome, and each panel corresponds to a different heterogeneity analysis. The variables Sales last month and Sales last 12 months are standardized with a mean of zero and a standard deviation of one. N = 109. Robust standard errors are in parentheses. P-values are $\leq 0.01^{***}$,  $\leq 0.05^{**}$, and $\leq 0.1^{*}$.
\end{minipage}
\end{tablenotes}

\end{threeparttable}
\end{table}

Column 1 of Table \ref{tab:hetero} presents results using sales as the outcome variable, revealing strong and statistically significant heterogeneity in treatment effects. While treated individuals with average prior sales experience a positive effect of 25.32 units (p < 0.01), the interaction term between treatment and lagged sales is also positive and highly significant (53.52, p < 0.01), indicating that the effect of the intervention increases with individuals’ sales performance in the previous month. Specifically, the estimated treatment effect rises to approximately 79 units for those one standard deviation above the mean in prior sales, whereas it becomes negative for those one standard deviation below the mean. Turning to the additional outcomes in Columns 2–6, the results are somewhat mixed. While the treatment effect on the log-transformed sales variable (Column 2) is positive but statistically insignificant, the pattern of heterogeneity remains in the expected direction, though the interaction term is not statistically significant. For electronic sales (Column 3), the treatment effect is large and significant (22.46, p < 0.01), with a strong and significant interaction term (38.58, p < 0.01), mirroring the findings from Column 1 and suggesting that the overall sales effects are largely driven by increases in electronic transactions. No significant treatment effect is found for cash sales (Column 4), and the interaction term is also insignificant, implying that the intervention did not affect cash-based transactions, which is reasonable since the intervention did not target them. In terms of labor supply, the treatment significantly increases both hours worked (Column 5) and number of working days (Column 6), with positive and significant interaction terms in Column 5 but not in Column 6, suggesting some heterogeneity in effort responses, particularly in terms of hours worked.

Overall, irrespective of how experience is measured (Panels a–c), the results are consistent: treatment effects are stronger for more experienced sellers than for less experienced ones, which aligns with previous research. Experienced sellers sold more magazines, worked more hours, and potentially more days, whereas inexperienced sellers were unaffected, or even decreased their sales, due to the treatment. However, when considering age (Panel d), we observe few differences in response to the treatment, implying that there is no clear heterogeneity corresponding to the age of the sellers. A priori, it is not obvious that such differences should exist.

It should be noted that distinguishing experience from general interest in selling many papers is difficult because there is likely some survivorship bias in the sample since only individuals with a relative fondness for selling papers are likely to do so for a long time and receive their main source of income from this activity. Thus, our results on experience are consistent with experienced sellers responding more strongly to greater earnings potential, but the results may also be due to, for instance, stronger responses from those for whom the paper is a more important source of income. Interpreting experience in the taxi studies would seem to potentially suffer from the same problem. 

\subsection{Robustness}
\textit{No control variables:} As our first robustness check, we estimate the regressions from equation~\eqref{eq:regression} without control variables. As expected, given the imbalance between treated and controls found in Appendix Table \ref{tab:balance} and lower power when not including lags with predictive power, all effect estimates are smaller and imprecisely measured in models without controls for past selling behavior, see panel a in Appendix Table \ref{tab:robustness}.

\textit{Full set of controls:} Secondly, we estimate the regressions from equation~\eqref{eq:regression} with the full set of control variables. This controls for the imbalance between treated and waitlist control, but is costly considering our small sample size, see panel b in Table \ref{tab:robustness}. While the results are similar to those in Table \ref{tab:main}, the estimate for Sales attenuates substantially and loses statistical significance, whereas most other outcomes show smaller changes. Accordingly, inference is weaker in this specification (standard errors are not uniformly larger across outcomes).

\textit{Low-dimensional ANCOVA benchmark:} Finally, to assess whether our main results are sensitive to the use of machine-learning–based covariate selection, Table~\ref{tab:sensitivity} reports transparent low-dimensional benchmarks. Column (1) presents an ANCOVA specification controlling only for the pre-treatment average of sales, yielding a positive and statistically significant treatment effect of 32.76 (p < 0.05). Column (2) shows that a parsimonious model controlling for a single one-period lag of sales (Sales$_{12}$) and basic demographics also delivers a positive and significant estimate of 25.21 (p < 0.05). Column (3) reproduces the main post-double LASSO estimate (23.25, p < 0.05) for comparison. Across these specifications, the treatment effect remains positive and statistically significant, indicating that the estimated sales increase is not primarily driven by the covariate selection procedure but that controls for lagged sales are important.

\textit{Large share of electronic sales:} Following our pre-analysis plan, we estimate the effects of Hours Worked and Working days on a sample of sellers with a large share of electronic sales during the past three months (more than the median, which was 0.80). The reason for also using this sample to analyse the working time is that the working time measures should be more accurate for sellers who mainly sell via the electronic system. For someone selling a large fraction of their magazines for cash, working days and working time, as estimated by electronic sales, will not be accurate. Since someone with a large share of electronic sales may still have few records of sales if volumes are small, we also conducted the analyses on those selling more magazines than the median during the past three months (93 magazines). Estimates for Hours worked and Working Days were at least as pronounced as in the full sample, as shown in Appendix Table \ref{tab:swishshare}. Since these subsamples likely have more precise measures of working time, the results lend greater credibility to the findings in the full sample.

\textit{Effects on the waitlist control group:} Since the control group was treated after the treatment group (waitlist control), we can also analyze the effects of the treatment on the control group. In Appendix Table \ref{tab:waitlist}, as in the main analysis, we present results for the waitlist control group using equation~\eqref{eq:regression} and the post-double LASSO selection approach. However, here, the previously treated individuals serve as the control group, and the original treatment period (the issue released on September 25th) is excluded from the analysis. These estimates are imprecisely measured; however, the point estimates suggest small positive effects on sales. The lower bounds of the confidence intervals do not include negative effects of 20 percent, which would have implied full income targeting. One reason for being cautious in interpreting these findings is that there could potentially be longer-run effects of the intervention, and using the previously treated individuals as controls may be problematic. Additionally, \textit{Situation Sthlm} announced a price increase to the sellers on November 18th, raising the selling price from 80 to 100 SEK for the upcoming issue to be released on November 27th, and therefore the remuneration to the sellers from 40 to 50 SEK, in important ways making permanent the temporarily higher earnings from our experiment.

\textit{Analyses including both the treatment and waitlist control group:} We also estimate two--way fixed effects (TWFE) models to assess the effect of the intervention. These estimations suffer from the same problems with using the previous intervention group as controls, but it may still be interesting to estimate an average treatment effect for both groups jointly. We estimate the average treatment effect using
\begin{equation}
    Y_{it} = \alpha_i + \lambda_t + \beta T_{it} + \varepsilon_{it},
    \label{eq:twfe}
\end{equation}
where $Y_{it}$ is the outcome, and $T_{it}$ equals one if individual $i$ is treated during period $t$ and zero otherwise. Table \ref{tab:did_results} reports these results, pointing in the same direction as the main results.

To explore potential dynamics (and following the pre-analysis plan), we also estimate a corresponding event--study specification. All models include individual and period fixed effects, and standard errors are clustered at the individual level. Results are shown in Figure~\ref{fig:event_study}. The coefficients are normalized relative to the period immediately before treatment ($t=-1$), which is shown as a zero-valued reference point. Pre-treatment estimates are interpreted as evidence on pre-trends, while the estimates at $t=0$ and for post-treatment periods capture changes in sales after treatment relative to $t=-1$. The dashed vertical line marks the reference (baseline) period ($t=-1$). Confidence intervals are constructed using standard errors clustered at the individual level. The event-study results are qualitatively similar but less precisely estimated and only marginally significant.

\addtocounter{table}{-1}
\begin{table}[H]
\caption{\label{fig/event_study}Two-way fixed effects estimates}
\label{tab:did_results}
\centering

% Lägg hela threeparttable inuti adjustbox
\begin{adjustbox}{max width=\textwidth}
  \begin{threeparttable}

    \begingroup
    \small % endast tabellen blir mindre
    % VIKTIGT: did_results.tex ska börja med \begin{tabular} (inte \begin{table})
    \input{results/did_results.tex}
    \endgroup

    \begin{tablenotes}[flushleft] % samma bredd som treparttable (dvs tabellen)
    \small
    \item \textit{Notes:} Standard errors clustered at the individual level are in parentheses. 
    P-values are $\leq 0.01^{***}$,  $\leq 0.05^{**}$, and $\leq 0.1^{*}$.
    \end{tablenotes}

  \end{threeparttable}
\end{adjustbox}

\end{table}





\begin{figure}[H]
\centering
\includegraphics[width=\textwidth]{results/event_study.pdf}
\caption{Estimated effects on sales (event study design). 
Notes: The dashed line marks the reference period ($t=-1$). Vertical bars represent 95\% confidence intervals.
}
\label{fig:event_study}
\end{figure}


\subsection{Survey Evidence on Motivation and Perceptions}
To complement the experimental findings, we conducted a short post-experimental survey in December 2024 (see Appendix Section~\ref{sec:surv} for the translated questionnaire). The survey was not included in the pre-analysis plan and should be interpreted as exploratory. As an incentive, respondents received two free copies of the magazine to complete the survey. Of the 109 experiment participants, 63 completed the survey (58 percent), with a roughly even split between the treatment (31) and control (32) groups.

Appendix Table~\ref{tab:surveysummary} summarizes the responses. While 63 individuals returned the survey, the effective sample size for each item is slightly smaller due to item non-response. A majority of sellers (52 percent) report having an earnings target, and 68 percent believe that working more leads to higher sales. However, only ten percent say they worked more during the bonus period, despite experimental evidence indicating significant increases in hours worked and magazines sold. This discrepancy between stated motivations and revealed behavior suggests that self-perceptions may not align with actual effort responses.

Stated preferences around work structure and income smoothness vary across individuals. Half (50 percent) report working in a single stretch during the day, and nearly one-third (31 percent) say they would prefer a lower, fixed wage over a higher, more variable one. Most respondents (66 percent) claim that higher income in one month does not lead them to work less in the next, and only 25 percent report adjusting working hours based on expected earnings. The primary motivation for selling the magazine is financial for 63 percent of respondents.

Regarding time preferences and risk attitudes, 64 percent claim to prefer receiving 100 SEK immediately over 150 SEK in one month, indicating high discount rates, and 72 percent prefer the expected value to an uncertain lottery. These responses may help explain the general preference for immediate, stable income, though they do not appear to translate directly into income targeting behavior in the field experiment.

\subsection{Using Observational Data}
Our results contrast with some observational studies from the taxi industry that indicate income targeting. One possible reason is that cabdrivers may face common supply-side shocks. If they do not want to work on holidays, hours will be limited, and wages will be higher on those days. And high wages may cause cabdrivers to work a few hours on days when they did not plan to work. Another possible reason for the differences is that working hours are difficult to measure. To measure working time, the taxi studies use trip data, covering trips taken in sequence on a given day, along with the duration and fare of each trip. One potential shortcoming is that working time is overestimated on days when business is slow because occupying a car is not the same as working, since one may instead run personal errands. Once a cab is leased (the usual arrangement when not operating one’s own car) the cost is sunk. Previous literature is well aware of these potential difficulties, focusing on stopping decisions (\citet{agarwal2015singaporean}), or using instrumental variables to avoid mismeasurement of hours (\citet{camerer1997labor}. Nevertheless, estimates of elasticities using non-experimental approaches may be biased.

A naïve exemplification of such bias is shown in Table~\ref{tab:observational}. In the spirit of \citet{lalonde1986evaluating}, it shows estimates of income elasticities based on historical variations in hourly earnings for our sample of sellers. These findings reveal a pronounced contrast between experimental and non-experimental estimates. The outcome variable is the natural logarithm of daily hours worked, and the key explanatory variable is the natural logarithm of hourly earnings in the year preceding the intervention. Since working time should be more precisely measured on days with many sales, and for sellers who sell many magazines, we conduct our analysis on four different samples. Column~(1) includes all 109 sellers and all days on which they sold at least two magazines. Column~(2) restricts the sample to days with more than five magazines sold, while Column~(3) further restricts it to days with more than ten magazines sold. In Column~(4), we limit the sample to sellers whose total sales in the past year exceeded the median. Across all specifications based on observational data, we find consistently negative income elasticities, suggestive of strong, if not perfect, income targeting.

\begin{table} [H]
\centering
\begin{threeparttable} 
\caption{\label{tab:observational} Observational findings}
\small{\input{results/observational}}
\begin{tablenotes}
\textit{Notes:} This table reports fixed effects estimations. The dependent variable is the natural log of daily working hours, and the independent variable is the natural log of hourly earnings. Column~(1) includes all sellers and all days on which they made at least two sales. Column~(2) restricts the sample to days with more than five sales, Column~(3) restricts it to days with more than ten sales. Column~(4) limit the sample to sellers whose total sales in the past year exceeded the median. Robust standard errors are in parentheses. P-values are $\leq 0.01^{***}$,  $\leq 0.05^{**}$, and $\leq 0.1^{*}$. 
\end{tablenotes}
\end{threeparttable}
\end{table}


\section{Conclusion}
Treated sellers sold more papers, worked longer hours, and took fewer days off. Overall, we find no evidence of income targeting among sellers of the street paper \textit{Situation Sthlm}. Following a coin toss to randomly assign treatment status to sellers, our measures indicate that treated sellers work longer, earn more, and that the effect may be more concentrated among more experienced individuals.

\cite{fehrgoette} provides a particularly relevant point of comparison, since they also study labor-supply responses to a temporary increase in earnings in a randomized field experiment. In their setting—Swiss bicycle messengers paid on commission—a temporary increase in the commission rate led to a clear increase in labor supply. 
Our results are consistent with the central takeaway from that experiment: when workers have substantial discretion over labor supply and face a temporary, salient improvement in earning opportunities, overall work effort can increase rather than decrease. Building on this experimental benchmark, our contribution is to provide equally clean causal evidence in a very different labor market, where both measurement and identification are often more challenging.

Like \cite{fehrgoette}, we therefore rely on a randomized controlled trial rather than an observational or correlational design. Treatment status was determined by a coin toss and assigned to participants in front of them, and causal effects were estimated using prespecified analyses. Combined with sales data from both before and after the intervention, this design gives us confidence that our findings reflect true labor-supply responses to temporary wage changes. A further advantage of our setting is measurement. In many observational contexts, distinguishing time spent actively working from time spent taking breaks or running personal errands is difficult, which can complicate interpretation and potentially contribute to differences across studies. Our study mitigates this concern by relying on a clearly defined bonus window and objective, transaction-based measures of output—verified Swish sales—tied directly to seller behavior. The RCT structure, together with the revealed-preference nature of our data, therefore offers a clean and credible test of income targeting in labor supply decisions.

Further complicating interpretation, self-reports of motivation and behavior may also be unreliable. In our post-experimental survey, more than half of the respondents reported having income targets, and only a small minority (10 percent) said they worked more because of the bonus. Yet experimental data show increases in work effort among treated sellers. These discrepancies suggest that stated preferences may not reflect actual behavior, reinforcing the value of revealed-preference methods in labor-supply research. They also cast doubt on survey-based evidence for income targeting, where participants might report income targeting without it influencing their actions.

\pagebreak
\bibliography{ref}


\pagebreak

\section*{Appendix}

\global\long\def\thetable{A\arabic{table}}
\setcounter{table}{0}
\global\long\def\theequation{A\arabic{equation}}
\setcounter{equation}{0}
\global\long\def\thefigure{A\arabic{figure}}
\setcounter{figure}{0}
\setcounter{section}{0}
\renewcommand{\thesection}{A.\arabic{section}}

% Ny rubrik A.1
\section{Tables}\label{app:tables}
% (lägg dina tabeller här, t.ex. \input{appendix_tables.tex})

\begin{table}[H]
\centering
\begin{threeparttable}
\caption{\label{tab:summary}Summary Statistics}
\begin{tabular}{lcccc}
\toprule
 & Mean & Standard Deviation & Minimum & Maximum \\
 & (1) & (2) & (3) & (4) \\
\midrule
Age & 55.83 & 10.87 & 26 & 82 \\
Woman & 0.34 & 0.48 & 0 & 1 \\
Sales & 112.85 & 191.07 & 0 & 1,555 \\
$\ln(\text{Sales})$ & 3.61 & 1.79 & 0 & 7.35 \\
Electronic & 96.36 & 177.08 & 0 & 1,467 \\
Cash & 16.50 & 25.38 & -10 & 160 \\
Hours & 30.01 & 42.68 & 0 & 286.7 \\
Days & 11.17 & 9.55 & 0 & 33 \\
\bottomrule
\end{tabular}
\begin{tablenotes}[flushleft]   % <-- gör noten vänsterjusterad
\small
\item \textit{Notes:} The number of observations for all variables is 109.
\end{tablenotes}
\end{threeparttable}
\end{table}
\clearpage



\begin{table}[H]
\centering
\begin{threeparttable}
\caption{\label{tab:balance} Balance}
\begin{tabular}{lcccc}
\toprule
 & Treatment Mean & Control Mean & Difference & P-value \\
\midrule
Age & 55.61 & 56.06 & 0.45 & 0.83 \\
Woman & 0.41 & 0.26 & -0.15 & 0.11 \\
Sales (t-1) & 78.55 & 93.23 & 14.67 & 0.62 \\
Sales (t-2) & 112.95 & 145.51 & 32.56 & 0.40 \\
Sales (t-3) & 67.73 & 92.17 & 24.44 & 0.37 \\
Sales (t-4) & 84.71 & 106.92 & 22.21 & 0.51 \\
Sales (t-5) & 61.98 & 99.77 & 37.79 & 0.18 \\
Sales (t-6) & 75.70 & 98.38 & 22.68 & 0.44 \\
Sales (t-7) & 51.73 & 94.94 & 43.21 & 0.03 \\
Sales (t-8) & 60.34 & 89.91 & 29.57 & 0.18 \\
Sales (t-9) & 135.04 & 155.87 & 20.83 & 0.69 \\
Sales (t-10) & 101.34 & 114.30 & 12.96 & 0.73 \\
Sales (t-11) & 76.04 & 97.75 & 21.72 & 0.47 \\
Sales (t-12) & 81.86 & 92.19 & 10.33 & 0.74 \\
\bottomrule
\end{tabular}
\begin{tablenotes}[flushleft]   % <-- nyckeln för vänsterjustering
\small
\item \textit{Notes:} Column 1 shows means of treated individuals ($n = 53$), Column 2 shows means of controls ($n = 56$). Column 3 reports the differences between the groups, and Column 4 reports p-values from t-tests of mean differences.
\end{tablenotes}
\end{threeparttable}
\end{table}
\clearpage



\begin{table}[H]
\centering
\begin{threeparttable}
\caption{\label{tab:correlations}Correlations between outcomes and lags}

\begin{tabular}{l*{6}{c}}
\toprule
            & Sales & ln(Sales) & Electronic & Cash & Hours & Days \\
\midrule
(t-1)       & 0.96  & 0.74      & 0.96       & 0.69 & 0.88  & 0.80 \\
(t-2)       & 0.83  & 0.70      & 0.80       & 0.49 & 0.76  & 0.75 \\
(t-3)       & 0.89  & 0.58      & 0.90       & 0.37 & 0.77  & 0.62 \\
(t-4)       & 0.89  & 0.54      & 0.90       & 0.36 & 0.76  & 0.53 \\
(t-5)       & 0.88  & 0.51      & 0.89       & 0.34 & 0.68  & 0.55 \\
(t-6)       & 0.89  & 0.57      & 0.89       & 0.43 & 0.69  & 0.57 \\
(t-7)       & 0.58  & 0.59      & 0.57       & 0.37 & 0.46  & 0.61 \\
(t-8)       & 0.79  & 0.61      & 0.81       & 0.29 & 0.67  & 0.57 \\
(t-9)       & 0.88  & 0.52      & 0.89       & 0.39 & 0.72  & 0.58 \\
(t-10)      & 0.88  & 0.54      & 0.89       & 0.41 & 0.73  & 0.57 \\
(t-11)      & 0.88  & 0.50      & 0.90       & 0.44 & 0.71  & 0.56 \\
(t-12)      & 0.87  & 0.44      & 0.88       & 0.39 & 0.72  & 0.55 \\
\bottomrule
\end{tabular}

\begin{tablenotes}[flushleft]
\small
\item \textit{Notes:} Each column shows the correlation of the outcome variable and the lags of the outcome variable.
\end{tablenotes}

\end{threeparttable}
\end{table}
\clearpage



\begin{table}[H]
\caption{\label{tab:robustness} No control variables, and full set of control variables}
\vspace{-0.3em}

\hfill
\begin{minipage}{0.9\textwidth}
\centering
\begin{threeparttable}

% --- Panel A ---
\begin{subtable}[t]{\linewidth}
\caption{\label{tab:main_nc} No control variables}
\centering
\begingroup
\small
\renewcommand{\arraystretch}{1.02}
\setlength{\abovetopsep}{0.35ex}
\setlength{\belowbottomsep}{0.35ex}
\setlength{\aboverulesep}{0.35ex}
\setlength{\belowrulesep}{0.35ex}
\input{results/main_nc}
\endgroup
\end{subtable}

\vspace{0.7em}

% --- Panel B ---
\begin{subtable}[t]{\linewidth}
\caption{\label{tab:main_fullcont} Full set of control variables}
\centering
\begingroup
\small
\renewcommand{\arraystretch}{1.02}
\setlength{\abovetopsep}{0.35ex}
\setlength{\belowbottomsep}{0.35ex}
\setlength{\aboverulesep}{0.35ex}
\setlength{\belowrulesep}{0.35ex}
\input{results/main_fullcont}
\endgroup
\end{subtable}

\vspace{0.6em}
\begin{tablenotes}[flushleft]
\small
\item \textit{Notes:} Each column corresponds to a different outcome variable, and each panel corresponds to a different heterogeneity analysis variable. N = 109. Robust standard errors are in parentheses. P-values are $\leq 0.01^{***}$,  $\leq 0.05^{**}$, and $\leq 0.1^{*}$.
\end{tablenotes}

\end{threeparttable}
\end{minipage}
\hfill

\end{table}


\addtocounter{table}{-1}

\begin{table}[H]
\centering
\caption{\label{tab:sensitivity} Sensitivity of the treatment effect to selection of controls}
\vspace{-0.3em}

\begin{minipage}{0.9\textwidth}
\centering
\begin{threeparttable}

\begingroup
\small
\renewcommand{\arraystretch}{0.9}
\setlength{\tabcolsep}{12pt}

% (valfritt) om du använder booktabs och vill styra luft runt linjer:
\setlength{\abovetopsep}{0.35ex}
\setlength{\belowbottomsep}{0.35ex}
\setlength{\aboverulesep}{0.35ex}
\setlength{\belowrulesep}{0.35ex}

\begin{tabular}{lccc}
\toprule
 & (1)  & (2)  & (3) \\
\midrule
Treated   & 32.76** & 25.21** & 23.25** \\
          & (14.23) & (10.60) & (9.33)  \\
Pre-period mean sales & 1.17*** &         &         \\
          & (0.10)  &         &         \\
Sales$_{12}$       &         & 1.21*** & 0.89*** \\
          &         & (0.06)  & (0.09)  \\
Age       &         & -0.33   &         \\
          &         & (0.58)  &         \\
Woman     &         & -23.50**&         \\
          &         & (10.91) &         \\
Sales$_{11}$ &      &         & 0.04    \\
          &         &         & (0.06)  \\
Sales$_7$ &         &         & 0.09    \\
          &         &         & (0.09)  \\
Sales$_4$ &         &         & 0.09    \\
          &         &         & (0.06)  \\
Sales$_3$ &         &         & 0.04    \\
          &         &         & (0.09)  \\
\midrule
Observations & 109 & 109 & 109 \\
\bottomrule
\end{tabular}

\endgroup

\vspace{0.2em}
\begin{tablenotes}[flushleft]
\small
\setstretch{1.5}
\textit{Notes:} Robust standard errors in parentheses.
Column (1) reports a transparent low-dimensional ANCOVA controlling for the pre-treatment average of sales.
Column (2) controls for one pre-period outcome (Sales$_{12}$) and demographics.
Column (3) reports an OLS regression controlling for the PDS-selected controls from the post-double LASSO procedure. P-values are $\leq 0.01^{***}$,  $\leq 0.05^{**}$, and $\leq 0.1^{*}$.
\end{tablenotes}

\end{threeparttable}
\end{minipage}
\end{table}





\begin{table}[H]
\caption{\label{tab:swishshare} Results for hours worked and working days}
\centering

\begin{adjustbox}{max width=\textwidth}
  \begin{threeparttable}

    \begingroup
    \small
    % Lite tajtare rader om du vill:
    \renewcommand{\arraystretch}{0.95}

    % l för första kolumnen (vänsterställd etikett),
    % sedan 4 fasta, centrerade kolumnbredder à 1.9 cm (justera vid behov)
    \begin{tabular}{l*{4}{>{\centering\arraybackslash}p{1.9cm}}}
    \toprule
     & (1) & (2) & (3) & (4) \\
     & Hours & Days & Hours & Days \\
    \midrule
    Treated & 19.52** & 4.09** & 19.32** & 2.82* \\
            & (6.45)  & (1.42) & (5.90)  & (1.42) \\
    \midrule
    Mean for controls & 37.82 & 13.85 & 49.83 & 19.96 \\
    N                 & 54    & 54    & 54    & 54 \\
    \bottomrule
    \end{tabular}

    \endgroup

    \begin{tablenotes}[flushleft]
    \small
    \item \textit{Notes:} This table reports treatment effects on hours worked and working days under for individuals with a large share of electronic sales (Columns 1 and 2), and for individuals with many electronic sales (Columns 3 and 4). Robust standard errors are in parentheses. P-values are $\leq 0.01^{***}$, $\leq 0.05^{**}$, and $\leq 0.1^{*}$.
    \end{tablenotes}

  \end{threeparttable}
\end{adjustbox}

\end{table}
\clearpage


\begin{table}[H]
\centering
\begin{threeparttable}
\caption{\label{tab:waitlist} Results for the waitlist control, with originally treated serving as controls}
\normalsize
\input{results/waitlist}

\vspace{0.5em} % lite extra luft mellan tabellen och noten

\begin{tablenotes}[flushleft]
\small
\item \textit{Notes:} This table reports estimates of treatment effects for the main outcome, sales, and five additional outcomes that report sales or working time. Estimations were made using the post-double LASSO selection approach with a full set of potential control variables. Robust standard errors are in parentheses. P-values are $\leq 0.01^{***}$, $\leq 0.05^{**}$, and $\leq 0.1^{*}$.
\end{tablenotes}
\end{threeparttable}
\end{table}
\clearpage


\begin{table} [H]
\centering
\begin{threeparttable} 
\small
\caption{\label{tab:surveysummary} Descriptive statistics from survey}
\small{\input{results/surveysummary1}}
\begin{tablenotes}
\small
\textit{Notes:} See Appendix Section \ref{sec:surv}  for the survey questions. 
\end{tablenotes}
\end{threeparttable}
\end{table}
\clearpage

% Första PDF-sidan: blir A.2 Questionnaire
\includepdf[
  pages=1,
  pagecommand={\section{Questionnaire}\label{sec:surv}\thispagestyle{plain}},
  scale=0.90
]{English.pdf}

% Resterande PDF-sidor utan extra rubrik
\includepdf[
  pages=2-,
  pagecommand={\thispagestyle{plain}},
  scale=0.95
]{English.pdf}

% Pre-analysis plan: rubrik på första PDF-sidan
\includepdf[
  pages=1,
  pagecommand={\section{Pre-analysis plan}\label{sec:pap}\thispagestyle{plain}},
  scale=0.95
]{pap_prospect.pdf}

% Resterande sidor utan extra rubrik
\includepdf[
  pages=2-,
  pagecommand={\thispagestyle{plain}},
  scale=0.95
]{pap_prospect.pdf}


\end{document}

